\documentclass[11pt,a4paper]{article}
% Préambule commun aux documents de PythonIPM/documentation.
% Compilation : tectonic <fichier>.tex
% tectonic compile avec XeTeX, qui lit l'UTF-8 nativement : surtout PAS inputenc ni fontenc,
% faits pour pdfLaTeX. Ce sont eux qui décodaient mal les caractères multi-octets — « cœur »
% ressortait en « cur », et les tirets cadratins en caractères parasites.
\usepackage{lmodern}
\usepackage[french]{babel}
\usepackage{microtype}
\usepackage[a4paper,margin=2.4cm]{geometry}
\usepackage{amsmath,amssymb}
\usepackage{booktabs}
\usepackage{longtable}
\usepackage{array}
\usepackage{ragged2e}
\usepackage{xcolor}
\usepackage{tcolorbox}
\usepackage{enumitem}
\usepackage{fancyhdr}
\usepackage[hidelinks]{hyperref}

\definecolor{encadre}{HTML}{F4F6F8}
\definecolor{bordure}{HTML}{4A6FA5}
\definecolor{alerte}{HTML}{B03A2E}

% Encadré « à retenir » : la lecture non technique du passage qui précède
\newtcolorbox{retenir}[1][]{
  colback=encadre, colframe=bordure, boxrule=0.6pt, arc=2pt,
  left=8pt, right=8pt, top=6pt, bottom=6pt,
  fonttitle=\bfseries, title={#1}
}

% Encadré d'avertissement : un point ouvert, un écart assumé
\newtcolorbox{attention}[1][]{
  colback=white, colframe=alerte, boxrule=0.6pt, arc=2pt,
  left=8pt, right=8pt, top=6pt, bottom=6pt,
  % le bandeau de titre a déjà la couleur d'alerte en fond : un titre `coltitle=alerte`
  % écrirait du rouge sur du rouge. On garde le blanc par défaut, comme l'encadré « retenir ».
  fonttitle=\bfseries, title={#1}
}

\newcommand{\col}[1]{\texttt{#1}}
% \ensuremath et non \textsubscript : \ipm est employé aussi bien dans le texte que dans les
% formules, et « M\textsubscript{0} » placé en mode mathématique est une erreur — silencieuse en
% PDF, mais qui fait échouer la conversion vers Word.
\newcommand{\ipm}{\ensuremath{M_{0}}}

\setlist[itemize]{itemsep=2pt, topsep=4pt}
\setlist[description]{style=nextline, leftmargin=1.2cm, itemsep=5pt}

% La source du document apparaît en tête de chaque page. Les documents RGPH la redéfinissent
% par \renewcommand{\source}{...} AVANT \begin{document}.
\newcommand{\source}{EHCVM 2021}

\pagestyle{fancy}
\fancyhf{}
\fancyhead[L]{\small\itshape IPM Côte d'Ivoire --- \source}
\fancyhead[R]{\small\thepage}
\renewcommand{\headrulewidth}{0.3pt}

\renewcommand{\arraystretch}{1.15}

\renewcommand{\source}{PND 2026-2030 / IPM RGPH 2021}

\title{\bfseries L'IPM calculé sur le RGPH 2021 au service du PND 2026-2030\\[4pt]
\large Analyse d'alignement}
\author{Pipeline \col{PythonIPM}}
\date{Septembre 2026}

\begin{document}
\maketitle

\begin{abstract}
\noindent Ce document confronte l'indice de pauvreté multidimensionnelle calculé sur le
Recensement général de la population et de l'habitation 2021 au Plan National de Développement
2026-2030. Il établit, cible par cible, ce que l'un permet de mesurer de l'autre, ce qu'il ne
permet pas, et l'usage auquel il se prête le mieux. Un document jumeau traite de l'IPM calculé sur
l'EHCVM 2021 ; un troisième les compare et conclut.
\end{abstract}

\begin{retenir}[La conclusion en cinq lignes]
\begin{itemize}
  \item L'IPM RGPH couvre \textbf{sept des vingt-deux cibles} du PND qu'un indice de pauvreté peut
        mesurer --- moins que l'enquête, faute de toute question de santé.
  \item Il est en revanche \textbf{le seul} à rendre opérationnels trois instruments centraux du
        plan : les \textbf{Pôles Économiques Compétitifs} du Pilier 5, le \textbf{registre social
        unifié} du Pilier 6 et le ciblage des \textbf{827 000 ménages} de filets sociaux du
        Pilier 4.
  \item Il publie un IPM pour \textbf{500 des 519 sous-préfectures et communes} et pour les
        \textbf{111 départements}, sans exception.
  \item Il établit que le milieu rural porte \textbf{84,7~\%} de la pauvreté
        multidimensionnelle du pays pour 45,8~\% de sa population.
  \item Étant décennal, il fournit \textbf{une photographie de départ}, pas un instrument de
        suivi : il ne dira pas en 2030 si le PND a produit ses effets.
\end{itemize}
\end{retenir}

\tableofcontents
\bigskip

\section{Ce que le PND demande à un indice de pauvreté}

Le PND 2026-2030 vise à « Bâtir une Grande Nation Stable, Ambitieuse et Solidaire ». Parmi ses
objectifs généraux figure de \textbf{« baisser significativement le taux de pauvreté »}, sans
valeur cible à 2030 --- seul objectif majeur à ne pas être chiffré.

Mais le plan ne demande pas seulement de \emph{mesurer}. Trois de ses instruments supposent de
savoir \emph{où} et \emph{qui}, à une maille que la statistique d'enquête n'atteint pas :

\begin{description}
  \item[Pilier 5] « renforcer la planification spatiale des investissements, fondée sur une
    \textbf{cartographie fine} des potentialités économiques régionales », développer les
    \textbf{Pôles Économiques Compétitifs} « à travers une répartition optimale et équitable » et
    aménager le territoire par les \textbf{villes secondaires} ;
  \item[Pilier 6] instaurer un \textbf{registre social unifié} « visant à renforcer l'inclusion et
    le ciblage des politiques publiques », et poursuivre « le renforcement du processus de
    décentralisation et de réduction des inégalités régionales » ;
  \item[Pilier 4] porter à \textbf{827 000} le nombre de ménages pauvres bénéficiaires de
    transferts monétaires au titre des filets sociaux.
\end{description}

\begin{retenir}[Deux questions, pas une]
Le PND pose au statisticien deux questions de nature différente. \emph{Où en est-on, et où en
sera-t-on en 2030 ?} --- c'est une question de mesure. \emph{Où faut-il investir, et quels ménages
faut-il servir ?} --- c'est une question de ciblage. Le recensement répond mal à la première et
très bien à la seconde.
\end{retenir}

\section{La correspondance, cible par cible}
\label{sec:correspondance}

Les treize indicateurs de l'IPM RGPH se rapportent aux cibles du PND comme suit. La colonne
« Aujourd'hui » donne la part de la population vivant dans un ménage privé sur l'indicateur.

\subsection{Pilier 4 --- Capital humain, compétences et emplois décents}

\begin{center}
\small
\begin{tabular}{@{}>{\RaggedRight}p{5.3cm}>{\RaggedRight}p{4.6cm}rc@{}}
\toprule
\textbf{Cible PND à 2030} & \textbf{Indicateur IPM} & \textbf{Auj.} & \textbf{Prise} \\
\midrule
Analphabétisme adultes 51,1 $\to$ 45~\% & Alphabétisation & 70,4~\% & directe \\
Analphabétisme 15-24 ans 37,6 $\to$ 20~\% & Alphabétisation & 70,4~\% & directe \\
Achèvement 1\textsuperscript{er} cycle secondaire
  81,14 $\to$ 89~\% & Année de scolarité (par le diplôme) & 63,2~\% & approchée \\
Chômage et main-d'œuvre potentielle (SU3) 12,9 $\to$ 8,9~\% & Chômage & 5,9~\% & directe \\
Espérance de vie 62,3 $\to$ 65 ans & Mortalité dans le ménage & 1,1~\% & indirecte \\
827 000 ménages bénéficiaires des filets sociaux &
  \textbf{identification exhaustive des ménages pauvres} & --- & \textbf{opérationnelle} \\
Enrôlés à la CMU 20 $\to$ 30 millions & --- & --- & \textbf{aucune} \\
Couverture par la protection sociale 27 $\to$ 50~\% & --- & --- & \textbf{aucune} \\
Personnes à moins de 5 km d'un centre de santé 82 $\to$ 100~\% & --- & --- & \textbf{aucune} \\
Retard de croissance des moins de 5 ans 21,4 $\to$ 18~\% & --- & --- & \textbf{aucune} \\
Indice de l'écart genre 0,655 $\to$ 0,7 & --- & --- & aucune \\
Préscolaire ; EFTP ; IPS primaire & --- & --- & aucune \\
\bottomrule
\end{tabular}
\end{center}

\subsection{Pilier 5 --- Infrastructures, pôles régionaux et transition écologique}

\begin{center}
\small
\begin{tabular}{@{}>{\RaggedRight}p{5.3cm}>{\RaggedRight}p{4.6cm}rc@{}}
\toprule
\textbf{Cible PND à 2030} & \textbf{Indicateur IPM} & \textbf{Auj.} & \textbf{Prise} \\
\midrule
Pôles économiques compétitifs, villes secondaires,
  réduction des inégalités régionales &
  \textbf{IPM à la maille de 519 sous-préfectures} & --- & \textbf{opérationnelle} \\
Accès à une source d'eau améliorée 81 $\to$ 90~\% & Eau potable (source seule) & 14,3~\% &
  directe \\
Accès à l'assainissement amélioré 37 $\to$ 70~\% & Toilette (sans le partage) & 45,1~\% &
  directe \\
Accès universel à l'électricité d'ici 2030, en priorité
  dans les quartiers précaires & Électricité & 11,5~\% & directe \\
Habitat décent (défi n\textsuperscript{o}~5 du diagnostic) & Logement & 25,0~\% & directe \\
Population utilisant internet 40,7 $\to$ 60~\% & Biens d'équipement & 18,2~\% & indirecte \\
Promiscuité, densité d'occupation du logement & --- & --- & aucune \\
\bottomrule
\end{tabular}
\end{center}

\subsection{Piliers 2, 3 et 6}

\begin{center}
\small
\begin{tabular}{@{}>{\RaggedRight}p{5.3cm}>{\RaggedRight}p{4.6cm}rc@{}}
\toprule
\textbf{Cible ou orientation PND} & \textbf{Indicateur IPM} & \textbf{Auj.} & \textbf{Prise} \\
\midrule
Registre social unifié pour le ciblage (P.~6) &
  \textbf{base exhaustive de 5,6 M de ménages} & --- & \textbf{opérationnelle} \\
Enregistrement des naissances 56,6 $\to$ 100~\% (P.~6) &
  Déclaration d'état civil & 15,2~\% & directe \\
NNI 27,78 $\to$ 100~\% (P.~6) & Déclaration d'état civil & 15,2~\% & indirecte \\
Modernisation de l'agriculture et productivité (P.~2) ;
  réduction de l'informalité (P.~3) &
  Emploi agricole de subsistance (approché) & 46,0~\% & approchée \\
Énergie de cuisson propre (P.~5) & Énergie de cuisson & 62,7~\% & directe \\
\bottomrule
\end{tabular}
\end{center}

\begin{retenir}[Le décompte]
Sur les vingt-deux cibles et orientations du PND qu'un indice de pauvreté peut éclairer, l'IPM
RGPH en couvre \textbf{sept de façon directe}, deux par approximation, deux de façon indirecte, et
en rend \textbf{trois opérationnelles} — le ciblage territorial, le registre social et la sélection
des ménages bénéficiaires — que l'enquête ne peut pas servir du tout. Huit restent hors de portée,
dont les quatre cibles santé et nutrition du Pilier 4.
\end{retenir}

\section{Ce que le RGPH est seul à apporter}

\subsection{Une carte publiable partout}

C'est l'apport décisif, et il se chiffre. Le tableau ci-dessous donne, pour chaque maille, le
nombre de modalités dont l'IPM est trop imprécis pour être publié --- convention usuelle : un
coefficient de variation supérieur à 20~\%.

\begin{center}
\begin{tabular}{@{}lrrr@{}}
\toprule
\textbf{Maille} & \textbf{Modalités} & \textbf{Non publiables} & \textbf{Population concernée} \\
\midrule
Ensemble & 1 & 0 & --- \\
Milieu de résidence & 3 & 0 & --- \\
Région & 33 & 0 & --- \\
Département & 111 & 0 & --- \\
Sous-préfecture ou commune & 519 & 19 & 1,55~\% \\
\bottomrule
\end{tabular}
\end{center}

\begin{retenir}[La maille du Pilier 5 est atteinte]
Le recensement publie un IPM pour \textbf{les 111 départements sans exception} et pour
\textbf{500 des 519 sous-préfectures et communes}. Les 19 restantes ne représentent que 1,55~\% de
la population, et leur imprécision tient au petit nombre de zones de dénombrement, non au nombre de
ménages.

C'est exactement l'échelle à laquelle le Pilier 5 entend conduire sa « répartition optimale et
équitable des investissements publics et privés ».
\end{retenir}

\subsection{Le diagnostic territorial que le PND appelle}

Le PND fait de la réduction des inégalités régionales une orientation du Pilier 6 et le
rééquilibrage territorial l'objet même du Pilier 5. L'IPM RGPH en donne la mesure.

\begin{center}
\begin{tabular}{@{}lrrrr@{}}
\toprule
\textbf{Milieu} & \textbf{Part pop.} & $H$ & \ipm & \textbf{Part de la pauvreté} \\
\midrule
Urbain & 44,8~\% & 5,4~\% & 0,0212 & 5,7~\% \\
Semi-urbain & 9,4~\% & 41,0~\% & 0,1705 & 9,7~\% \\
Rural & 45,8~\% & 71,1~\% & 0,3075 & \textbf{84,7~\%} \\
\bottomrule
\end{tabular}
\end{center}

Le rural et l'urbain pèsent le même poids démographique, et le rural porte quinze fois l'IPM de
l'urbain. À l'échelle des régions, l'écart va de 0,0069 (District autonome d'Abidjan) à 0,3654
(Bounkani), soit un rapport de \textbf{un à cinquante-trois}. À l'échelle des sous-préfectures, il
va de 0,0015 (Plateau) à 0,4752 (Santa de Ouaninou), où 92,6~\% de la population est
multidimensionnellement pauvre.

\begin{retenir}[Le semi-urbain, angle mort de l'enquête]
Le recensement distingue trois milieux là où l'EHCVM n'en code que deux. Le semi-urbain concentre
la population \textbf{vulnérable} : 34,2~\% de ses habitants ont un score de privation compris
entre 0,2 et le seuil de pauvreté, contre 17,7~\% en urbain et 20,9~\% en rural.

Ce sont précisément les « villes secondaires » que le Pilier 5 entend développer. L'enquête ne
permet pas de les isoler.
\end{retenir}

\subsection{Le ciblage : du territoire au ménage}

Le PND fixe un objectif de \textbf{827 000 ménages} bénéficiaires de transferts monétaires et
prévoit un \textbf{registre social unifié}. Ces deux instruments demandent une base exhaustive.

L'IPM RGPH identifie \textbf{1 751 746 ménages multidimensionnellement pauvres}, chacun avec son
score, son identifiant et sa localisation. L'objectif de 827 000 ménages représente donc
\textbf{47~\% des ménages pauvres} : le problème n'est pas d'en trouver assez, mais de choisir
lesquels --- ce qu'un score par ménage permet de faire par ordre de sévérité.

\begin{center}
\begin{tabular}{@{}lr@{}}
\toprule
\textbf{Statut} & \textbf{Ménages} \\
\midrule
Pauvreté sévère ($c_i \ge 0{,}5$) & 285 115 \\
Pauvres ($c_i \ge 1/3$) & 1 751 746 \\
Vulnérables ($0{,}2 \le c_i < 1/3$) & 1 155 195 \\
\bottomrule
\end{tabular}
\end{center}

Sur le plan géographique, \textbf{121 sous-préfectures concentrent la moitié} de la pauvreté
multidimensionnelle du pays, pour 41,5~\% de sa population, et 277 en concentrent 80~\%. Une
politique de ciblage géographique dispose là de sa liste de priorités.

\subsection{La mortalité, seule prise sur la santé}

Le recensement ne pose aucune question de santé, mais il recense les décès survenus dans le ménage
au cours des douze mois précédents. L'IPM en tire son unique indicateur de santé : le décès d'un
enfant de moins de 18 ans, qui touche \textbf{1,1~\%} des ménages.

C'est la seule prise, même indirecte, sur la cible « espérance de vie de 62,3 à 65 ans ».

\section{Ce que l'IPM RGPH ne peut pas faire pour le PND}

\subsection{La dimension Santé du Pilier 4 lui échappe entièrement}

\begin{attention}[Quatre cibles chiffrées sans aucune prise]
Le recensement ne permet de mesurer \textbf{aucune} des quatre cibles santé et nutrition du
Pilier~4 : les 30 millions d'enrôlés à la CMU, la couverture par la protection sociale portée à
50~\%, la population à moins de 5 km d'un établissement de santé portée à 100~\%, et la prévalence
du retard de croissance ramenée à 18~\%.

Ce n'est pas une lacune de l'indice, c'est une propriété du questionnaire : un recensement
administré à trente millions de personnes ne peut pas porter de module de santé.
\end{attention}

La conséquence se lit dans la décomposition de l'indice. La dimension Santé, réduite à la seule
mortalité, porte un quart du poids nominal et ne contribue qu'à \textbf{1,6~\%} de \ipm{}.

\begin{center}
\begin{tabular}{@{}lrr@{}}
\toprule
\textbf{Dimension} & \textbf{Poids nominal} & \textbf{Contribution à \ipm} \\
\midrule
Éducation & 25~\% & 40,6~\% \\
Conditions de vie & 25~\% & 29,2~\% \\
Emploi & 25~\% & 28,6~\% \\
Santé & 25~\% & \textbf{1,6~\%} \\
\bottomrule
\end{tabular}
\end{center}

\subsection{Le suivi dans le temps}

\begin{attention}[Une photographie, pas un film]
Le RGPH est décennal. Le prochain n'interviendra pas avant la fin du PND 2026-2030. L'IPM RGPH
fournit donc une \textbf{valeur de départ} très solide et très détaillée, mais ne pourra pas dire,
à mi-parcours ni en 2030, si les cibles ont été atteintes.

Le PND prévoit un cadre de gouvernance chargé d'apprécier les progrès : le recensement ne peut pas
l'alimenter en évolution.
\end{attention}

\subsection{Trois autres manques}

\begin{description}
  \item[Le genre.] Le sexe du chef de ménage n'est pas dans l'extrait transmis. La cible « indice
    de l'écart genre 0,655 $\to$ 0,7 » ne peut pas être éclairée par cet indice.
  \item[La pauvreté monétaire.] Le recensement n'a pas de module de consommation. L'objectif
    « baisser significativement le taux de pauvreté », que le PND documente sur la série monétaire,
    ne peut pas être croisé avec la pauvreté multidimensionnelle dans cette source.
  \item[La promiscuité.] Le nombre de pièces du logement n'est pas collecté, alors que le
    développement urbain et l'accès à un logement décent figurent parmi les dix défis du diagnostic
    stratégique du PND.
\end{description}

\subsection{Deux approximations à signaler}

Deux indicateurs reposent sur une substitution qu'il faut mentionner si le chiffre est publié dans
un document du PND.

\begin{description}
  \item[L'année de scolarité] est mesurée par le \emph{diplôme} le plus élevé --- le BEPC ou plus
    vaut dix années accomplies --- faute de question sur les années d'études ou la classe atteinte.
    L'approximation est bonne au niveau national (63,2~\% contre 62,9~\% à l'enquête), moins bonne
    individuellement.
  \item[L'emploi agricole de subsistance] retient le chef de ménage occupé dans une branche
    agricole, à son compte ou comme aide familial, faute de question sur « son propre champ ».
    C'est le \textbf{premier contributeur à l'indice} (26,8~\%) : l'approximation pèse lourd.
\end{description}

\section{L'usage auquel l'IPM RGPH se prête}

\begin{retenir}[Trois emplois immédiats]
\begin{enumerate}
  \item \textbf{La carte de référence du Pilier 5.} Un IPM publiable pour 111 départements et 500
        sous-préfectures, qui donne à la « planification spatiale des investissements » sa base
        statistique.
  \item \textbf{Le socle du registre social unifié du Pilier 6.} 5,6 millions de ménages, chacun
        avec un score de privation, un identifiant et une localisation — dont 1,75 million de
        pauvres et 285 000 en pauvreté sévère.
  \item \textbf{La liste de priorités des filets sociaux.} 121 sous-préfectures concentrent la
        moitié de la pauvreté du pays ; l'objectif de 827 000 ménages peut y être décliné par ordre
        de sévérité.
\end{enumerate}
\end{retenir}

Il ne se prête en revanche pas au cadre de résultats du plan : sept cibles seulement, aucune sur la
santé, et aucune possibilité d'actualisation avant 2030.

\section{Limites de cette analyse}

\begin{enumerate}
  \item \textbf{Les correspondances ne sont pas des identités.} Les indicateurs de l'IPM sont des
        indicateurs de \emph{ménage} ; les cibles du PND sont le plus souvent des taux
        \emph{individuels}. Les deux évoluent dans le même sens sans être le même nombre.
  \item \textbf{L'IPM RGPH n'est pas comparable en niveau à l'IPM EHCVM}, la dimension Santé du
        premier étant presque inerte. Toute comparaison entre les deux sources doit passer par la
        variante harmonisée.
  \item \textbf{Un écart de taux de privation sur la toilette} (45,1~\% ici, 33,09~\% dans la
        production Stata antérieure) n'est pas réconcilié et reste à élucider avec l'INS.
  \item \textbf{Les intervalles de confiance d'un recensement} ne mesurent pas une incertitude
        d'échantillonnage mais une dispersion entre zones de dénombrement. Le seuil de publication
        à 20~\% de coefficient de variation reste une convention prudente, non une règle absolue.
  \item \textbf{Le PND ne contient pas de cible d'IPM.} L'alignement établi ici est un alignement de
        \emph{contenu}, pas une correspondance institutionnelle.
\end{enumerate}

\end{document}
